\chapter{FUNDAMENTAÇÃO TEÓRICA}
\label{cap:fundamentacao}

\section{\uppercase{Roteamento por Caminhos Mínimos e ECMP}}

Em uma rede IP, cada enlace recebe um peso (métrica) e os roteadores calculam, com o algoritmo de Dijkstra, os caminhos de menor peso total até cada destino. Quando existem vários caminhos mínimos entre dois nós, o mecanismo ECMP (\textit{Equal-Cost Multi-Path}) divide o tráfego entre eles.

Seja $G = (V, A)$ a rede, com métrica $\omega(a)$ em cada arco. O grafo de encaminhamento $FG(u, v)$ é o subgrafo formado por todos os arcos que pertencem a algum caminho mínimo de $u$ até $v$. Pela regra ECMP, o fluxo que chega a um nó de $FG(u, v)$ é dividido igualmente entre os arcos de saída desse nó em $FG(u, v)$. A fração do fluxo de $u$ para $v$ que passa pelo arco $a$ é o coeficiente de divisão $r(u, v, a)$, nulo se $a \notin FG(u, v)$.

\section{\uppercase{Segment Routing}}

O \textit{Segment Routing} (SR) permite que a origem defina uma sequência de nós que o pacote deve visitar, codificada no próprio cabeçalho. Um caminho de segmentos para uma demanda de $s$ para $t$ é uma sequência $\langle s, w_1, \dots, w_\ell, t \rangle$, em que $w_1, \dots, w_\ell$ são os \textit{waypoints}. Cada par consecutivo é um segmento, e entre os extremos de um segmento o tráfego segue o grafo de encaminhamento $FG$ do IGP. Sem \textit{waypoints} ($\ell = 0$), a demanda é roteada pelos caminhos mínimos de $s$ a $t$.

A vantagem do SR para a engenharia de tráfego é que poucos segmentos bastam para desviar demandas de enlaces congestionados, sem alterar a métrica de nenhum enlace e, portanto, sem afetar o restante do tráfego.

\section{\uppercase{Engenharia de Tráfego e Utilização dos Enlaces}}

A carga (utilização) de um enlace é a razão entre o fluxo que passa por ele e sua capacidade; valores acima de 1 indicam congestionamento. A utilização máxima dos enlaces (MLU) é o critério mais usado em engenharia de tráfego. Abordagens clássicas otimizam os pesos do IGP; com SR, a otimização passa a ser a escolha dos \textit{waypoints} de cada demanda. O desafio EURO/ROADEF 2026 estende o critério para uma ordem lexicográfica: minimiza-se o enlace mais carregado, depois o segundo mais carregado, e assim por diante, considerando todos os períodos.

\section{\uppercase{Meta-heurísticas de Busca em Vizinhança}}

Uma busca local parte de uma solução e a substitui repetidamente por uma solução vizinha melhor, até atingir um ótimo local. A vizinhança é definida por um tipo de movimento, como inserir ou trocar um \textit{waypoint}.

A \textbf{Descida em Vizinhança Variável} (VND) explora várias vizinhanças em ordem fixa: ao encontrar uma melhora, volta para a primeira; ao esgotar todas, termina em uma solução que é ótimo local para todas elas. A \textbf{Busca em Vizinhança Variável} (VNS) combina essa descida com perturbações (\textit{shaking}) que tiram a busca do ótimo local, aceitando a nova solução apenas se ela for melhor que a melhor conhecida.

Buscas adaptativas, como a ALNS (\textit{Adaptive Large Neighborhood Search}), sorteiam a vizinhança a aplicar em cada iteração, com pesos que aumentam para os movimentos que mais trouxeram melhorias. A AMNS usada neste trabalho segue essa ideia: os pesos dos movimentos se adaptam ao longo da busca, e a VNS é usada para escapar de ótimos locais.

\section{\uppercase{Trabalhos Relacionados}}

Fortz e Thorup (2000) mostraram que otimizar os pesos do OSPF reduz significativamente o congestionamento, mas qualquer mudança de peso afeta todo o tráfego que usa o enlace. Bhatia et al. (2015) formularam a engenharia de tráfego com \textit{Segment Routing} e mostraram que poucos segmentos por demanda já aproximam o roteamento ótimo. O desafio EURO/ROADEF 2026 introduz a versão temporal com intervenções e orçamento de reconfiguração; as soluções publicadas na fase \textit{sprint} do desafio são uma referência de comparação natural.

\citet{yang2023intelligent} tratam o balanceamento de carga com SR sob um limite para o tamanho do cabeçalho de segmentos, que representa o custo de controle do SR. No esquema proposto, o IRTE, um controlador centralizado coleta o estado da rede e decide, para cada fluxo, se ele segue o caminho mínimo ou um caminho de desvio que passa por nós de desvio escolhidos, com um limite para o número de rótulos de desvio por fluxo. Os autores formulam a escolha como um problema de mapeamento de fluxos a nós de desvio, mostram que ele é NP-difícil por redução ao problema do subgrafo $k$-denso e propõem dois algoritmos: o LBA-DR, baseado em programação linear com arredondamento aleatório, e o IACO, uma otimização por colônia de formigas. Em seis topologias reais do \textit{Internet Topology Zoo}, com 11 a 125 nós, a utilização máxima de banda foi até 24,6\% menor que a do roteamento tradicional, na topologia BellCanada. O objetivo (a carga do enlace mais carregado) e o limite de segmentos por caminho são os mesmos do problema tratado neste trabalho; o problema deles, porém, considera uma única configuração da rede, sem falhas nem custo de reconfiguração ao longo do tempo, e as redes avaliadas são bem menores que as do desafio, que chegam a 3902 nós.

\citet{aubry2018robustly} exploram uma propriedade do SR que também é central neste trabalho: quando um enlace falha, o IGP recalcula os caminhos mínimos e a mesma lista de segmentos passa a corresponder a outro caminho, sem nenhuma reconfiguração. Os autores definem pares de caminhos robustamente disjuntos, que permanecem disjuntos para qualquer falha de um conjunto dado pelo operador, e buscam o par com o menor atraso no pior caminho. Eles provam que o problema é NP-difícil em geral, mas que se torna tratável em tempo polinomial quando o número máximo de segmentos por caminho é fixo, e usam uma condição necessária e suficiente para a robustez para podar o espaço de busca. Em topologias do Rocketfuel e do \textit{Internet Topology Zoo}, os caminhos são calculados em frações de segundo ou poucos segundos na maioria dos casos, e um único nó intermediário eleva a fração de pares origem-destino com caminhos robustamente disjuntos para 71\% a 100\%, enquanto desvios adicionais trazem pouco ganho. A fase robusta deste trabalho segue a mesma ideia, com outro objetivo: escolher, no primeiro período, listas de segmentos cujos caminhos continuem bons depois das intervenções dos períodos seguintes, recalculados pelo IGP. A pouca utilidade de desvios adicionais também tem um análogo aqui, com o teste da inserção dupla de \textit{waypoints} (Seção~\ref{sec:insercao_dupla}).

\citet{chen2022waypoint} aplicam SR a redes de satélites de órbita baixa (LEO), cuja topologia muda com o movimento dos satélites. O tempo é dividido em intervalos de topologia fixa, e os enlaces entre planos orbitais são desligados quando os satélites entram nas regiões polares. Como essas mudanças são previsíveis, os autores escolhem até dois \textit{waypoints} por caminho para que a rota evite, já no intervalo atual, os enlaces que serão desligados no intervalo seguinte. O problema é formulado como programação não linear inteira mista e resolvido em etapas: os enlaces que serão desligados são removidos do grafo, o caminho mínimo é calculado com o algoritmo de Dijkstra, e os \textit{waypoints} são escolhidos entre os nós da mesma órbita da origem e do destino. Em simulações com o NS-3, o algoritmo teve menos perda de pacotes, atraso e variação de atraso que o OLSR. Assim como no desafio, as indisponibilidades de enlaces são conhecidas com antecedência e as rotas são escolhidas para atravessá-las sem reconfiguração; o objetivo, porém, é a estabilidade e o número de saltos de cada caminho, e não a carga dos enlaces, e não há orçamento de reconfiguração.

Em resumo, os trabalhos revisados tratam separadamente aspectos que o problema T-ASR combina: o balanceamento de carga com poucos segmentos por caminho \citep{yang2023intelligent}, a robustez de listas de segmentos a falhas sem reconfiguração \citep{aubry2018robustly} e o planejamento de rotas para mudanças previsíveis na topologia \citep{chen2022waypoint}. Nenhum deles considera, ao mesmo tempo, a minimização da carga máxima em vários períodos com intervenções programadas e um orçamento que limita as mudanças de segmentos entre períodos, que é a combinação abordada neste trabalho.
